\section{The native graph}
\label{sec:native-graph}

We write $G60$ for the exact labeled sixty-vertex graph used throughout
this paper. The graph is identified with $\mathrm{AT4val}[60,6]$.

\subsection{Basic parameters}

The native graph has
\[
|V(G60)|=60,
\qquad
|E(G60)|=120,
\]
and every vertex has degree $4$. Thus $G60$ is quartic. The graph is
connected.

The automorphism calculations in later sections are performed on the
exact native edge set rather than on a graph reconstructed only from
numerical invariants. This distinction is important: matching order,
degree, spectrum, or other fingerprints does not by itself identify a
labeled graph or its permutation action.

\subsection{The canonical deck involution}

The native model carries a fixed-point-free involution
\[
a\colon V(G60)\longrightarrow V(G60).
\]
Its orbits are thirty disjoint pairs
\[
[v]_a=\{v,a(v)\}.
\]
We refer to these pairs as the \emph{$a$-fibers}. Since $a^2=1$ and
$a$ has no fixed vertices, the orbit set
\[
V(G60)/\langle a\rangle
\]
has cardinality $30$.

The involution $a$ is an automorphism of $G60$. Consequently, adjacency
descends to the orbit set and defines a two-fold quotient graph
\[
G60\longrightarrow G60/\langle a\rangle.
\]
We denote this quotient by $G30$.

The notation $a$ is retained because it is the native deck label in the
project data and certificates. When discussing general covering theory,
we may also describe $a$ as the canonical deck involution.

\subsection{Exact labeled data}

For a finite labeled graph, an automorphism is a permutation
\[
g\in\operatorname{Sym}(V(G60))
\]
such that
\[
\{u,v\}\in E(G60)
\quad\Longleftrightarrow\quad
\{g(u),g(v)\}\in E(G60).
\]

All permutation calculations in this paper use a fixed ordering of the
sixty native vertices and the corresponding exact edge list. A
permutation is therefore represented as a length-$60$ array, and group
multiplication is permutation composition in that fixed labeling.

The explicit labeling is not mathematically canonical, but it is
essential for reproducibility. The abstract conclusions are independent
of the chosen labels, while the exported certificates record the precise
labeling under which the computations were performed.

\subsection{Distinction from the thirty-vertex companion graph}

The quotient $G30$ is not isomorphic to the separately studied
thirty-vertex graph
\[
X=\operatorname{KC}(L(P)).
\]
Indeed,
\[
\left|\operatorname{Aut}(G30)\right|=240,
\qquad
\left|\operatorname{Aut}(X)\right|=720,
\]
so the two graphs cannot be isomorphic. The present manuscript concerns
the native graph $G60$ and the specific quotient arising from its
involution $a$. More generally, comparisons between thirty-vertex models
must be established by explicit invariants, isomorphisms, or certificates
rather than inferred from notation or vertex count.

In the next section we study the quotient map
\[
G60\longrightarrow G30
\]
and determine how quotient automorphisms lift to the native graph.
